% -------------------------------------------------
% Vstupní slovo pro širší automat
% -------------------------------------------------
\newcommand{\FAInputWide}[3]{%
    \node[font=\large] at (3.2,1.9) {$#1\textcolor{red}{#2}#3$};
}

\newcommand{\FAInputWideFinished}[1]{%
    \node[font=\large] at (3.2,1.9) {$#1\textcolor{red}{\lambda}$};
}

% -------------------------------------------------
% Jeden slide animace pro automat rozpoznávající podslovo 101
% #1 ... styl q0
% #2 ... styl q1
% #3 ... styl q2
% #4 ... styl q3
% #5 ... kód pro vstupní slovo
% #6 ... kód stromu výpočtu (+ případně výsledek)
% -------------------------------------------------
% -------------------------------------------------
% Jeden slide animace pro automat rozpoznávající podslovo 101
% #1 ... styl q0
% #2 ... styl q1
% #3 ... styl q2
% #4 ... styl q3
% #5 ... kód pro vstupní slovo
% #6 ... kód stromu výpočtu (+ případně výsledek)
% -------------------------------------------------
\newcommand{\SubstringAutomatonFrame}[6]{%
\begin{frame}{Výpočet deterministického automatu}{Příklad}
% Obrázek se zmenší tak, aby se vešel na snímek ve všech poměrech stran.
\begin{adjustbox}{max width=\linewidth, max totalheight=0.78\textheight, center}
\begin{tikzpicture}[>=stealth]

    % vstup
    #5

    % automat
    \initialstate[#1]{q0}{(0,0)}{$q_0$}{(-1.2,0)}{west}
    \state[#2]{q1}{(2.1,0)}{$q_1$}
    \state[#3]{q2}{(4.2,0)}{$q_2$}
    \acceptingstate[#4]{q3}{(6.3,0)}{$q_3$}

    % přechody
    \transitionarrow[loop above]{q0}{above}{$0$}{q0}
    \transitionarrow{q0}{above}{$1$}{q1}
    \transitionarrow[loop above]{q1}{above}{$1$}{q1}
    \transitionarrow{q1}{above}{$0$}{q2}
    \transitionarrow[out=130,in=50,looseness=1.5]{q2}{above}{$0$}{q0}
    \transitionarrow{q2}{above}{$1$}{q3}
    \transitionarrow[loop above]{q3}{above}{$0,1$}{q3}

    % strom výpočtu
    \node[anchor=west,font=\normalsize] at (-0.3,-1.45) {Strom výpočtu:};
    #6

\end{tikzpicture}
\end{adjustbox}
\end{frame}
}